% Appendix F: Background (promised to reviewer Fbco: "a Background section (LSTV subtypes;
% Castellvi I to IV; the VERIDAH architecture)"). Integrated 2026-09-27 (third pass) from the
% verified draft scratchpad/neurips_promises/drafts/app_background.tex; notes in
% app_background_NOTES.md. Uses sec:dataset, sec:castellvi, tab:castellvi_dist, sec:audit,
% sec:unmerged, app:spineps and app:tsv3 (appendix/spineps.tex). If the SPINEPS section is cut,
% replace "(Appendix~\ref{app:spineps})" with "(the current release)"; if the total_v3 slot is
% cut, end the footnote with "It was not evaluated here."
%
% Every factual sentence is sourced. Where each number was checked (27 Sep 2026):
%  Konin & Walz 2010, AJNR 31:1778: 4-30%; Castellvi definitions; morphology ("varying degrees
%    of lumbarization, such as ..." - the S1 features are examples, not all required); p.1778
%    (first page, not 1779; checked page-by-page in MANUAL_REFS/KONIN_WALZ.pdf) "does not provide
%    information relevant to accurate enumeration of the involved segment".
%  Castellvi 1984: classification; 60 of 200 myelography patients (PubMed 6495013 abstract,
%    saved in paper/mpda/_schehr/sources.json). Full text NOT read; the type definitions are
%    Konin & Walz's restatement and are attributed to them.
%  Carrino 2011: 147 subjects; LSTV associated with anomalous total count, OR 7.4 (Europe PMC
%    abstract). Sekharappa 2014: 3 x 1,000; 8.1/14/16.9%; type I "not called a transitional
%    vertebra in our study" (PMC3939369). Hanhivaara 2022: 310 of 1,468 (21.1%), Castellvi
%    classification used, type I findings reported (Europe PMC abstract). Nagata 2025: 551
%    whole-spine CTs, 16.3%; sacralization = 23 presacral, lumbarization = 25 (PubMed abstract).
%    Lian 2018: counting caudally from C2 is the gold standard (abstract).
%  Liebl 2021 (VerSe'20, PMC8553749): Wigh's definition (ribs < 38 mm on both sides = TLTV) is
%    quoted, but VerSe'20 itself labels a vertebra thoracic if at least one rib > 38 mm (with
%    diagonal downward alignment) and LUMBAR if both are shorter; Castellvi II-IV graded ("only
%    Castellvi II-IV were classified"); "did not segment vertebrae that showed partial fusion
%    with the sacrum (Castellvi grades III and IV)"; variants per subject in Suppl. Table 1.
%  VerSe ids: github.com/anjany/verse README ("20-25: lumbar spine: L1-L6", "28: ... T13").
%  CTSpine1K (Deng et al.): "Every vertebra (from C1 to L6) is labeled with integer values from
%    1 to 25"; "containing L6 (resulting from sacral lumbarization) or losing L5 (resulting from
%    lumbar sacralization)"; "the DSC of L6 is almost 0 ... the existence of L6 confuses the
%    model, resulting in prediction dislocations".
%  VERIDAH (arXiv:2601.14066v1, the only version): Sec. 2.1-2.2, Tables 2, 3, 5, Figs 1D, 5.
%    The label head is stated as "23 classes" while the cost set C is "all 24 vertebra labels"
%    C1..L5 (the paper's own inconsistency), so the text says only "no T13 or L6 class". The CT
%    cohort is "in-house" (not called private). Locations come from instance segmentations; the
%    method "requires the centroids of vertebrae beforehand".
%  SPINEPS paper (Eur Radiol 35:1178): MRI only; "the instance labels are counted from top to
%    bottom regardless of the field of view or the presence of enumeration abnormalities".
%  SPINEPS 2.1.2 (PyPI, uploaded 24 Sep 2026): CT_semantic, CT_instance, ct_labeling
%    (DenseNet169); post-map "L1..L6 -> 20..25, T13 -> 28" (scout_spineps/s212/.../phase_labeling.py).
%  TotalSegmentator: 'total' map = C1..L5, sacrum, S1, no L6/T13 (map_to_binary.py in the
%    2.4.0 wheel, 117 classes; README "Class details" table); vertebrae_S1 new in v2
%    (improvements_in_v2.md); the audit ran TotalSegmentator==2.4.0
%    (docker/Dockerfile.totalsegmentator @ c53d2dc). Footnote, re-checked 27 Sep 2026: the
%    PyPI wheels 2.15.0 (1 Jul 2026), 2.17.0 and 2.18.0 (12 Aug) all define
%    class_map["total_v3"] = total with [26]="vertebrae_L6" (2.12.0-2.14.0 do not);
%    map_tasks_config.py maps total_v3 to task ids 831-837 with version "v3.0.0-weights"; the
%    GitHub release v3.0.0-weights (pre-release, published 2026-09-07, body "Add weights for
%    v3.0.0") carries Dataset831-837 zips, and a range request for 836 returned HTTP 206, so the
%    weights ARE downloadable (not license-gated: total_v3 is not in commercial_models). Only the
%    "download all" script download_pretrained_weights.py omits 831-837. CHANGELOG and README
%    do not mention total_v3 or L6.
%  MDPA (arXiv:2609.22760v1, describes release v11): "no scan in this dataset contains C2";
%    v11 adds lower thoracic levels, ribs, lumbar ribs and classes for L6 and T13. April has 1,153
%    records, not v11's 802: the 351 extra series are the other acquisition of the same CT
%    colonography exams (ACRIN 6664), so "no C2" carries over by field of view, not by identity
%    of series. (v11 Table S3 lists id 1 "C1" in 2 records and id 6 "C6" in 1 record, isolated
%    labels with no C2-C5 or T1-T2 anywhere; the highest contiguous level is T3. Not checked
%    further; it does not bear on this appendix.)

\section{Background: transitional vertebrae, Castellvi grading and \VERIDAH{}}
\label{app:background}

\paragraph{Enumeration and transition.}
The spine usually has seven cervical, twelve thoracic and five lumbar vertebrae. An
enumeration anomaly gives eleven or thirteen thoracic or four or six lumbar vertebrae; a
transitional vertebra is a border segment with features of both neighbouring
regions~\citep{Moller2026VERIDAH:Sequences,Konin2010LumbosacralRelevance}. The two co-occur: a
lumbosacral transitional vertebra (LSTV) is associated with an anomalous total
count~\citep{Carrino2011EffectImaging}. An LSTV is either
sacralization of the lowest lumbar segment, from a broadened transverse process to complete
fusion with the sacrum, or lumbarization of the uppermost sacral segment, which may
articulate rather than fuse and may show lumbar-type facets or a fully sized S1--S2
disc~\citep{Konin2010LumbosacralRelevance}. The same words also name counts:
\citet{Nagata2025LSTVWholeSpine} define sacralization as 23 presacral vertebrae and
lumbarization as 25. Rib conventions differ as well: \VERIDAH{} calls any vertebra with a true
rib thoracic, whatever the rib's length~\citep{Moller2026VERIDAH:Sequences}; VerSe'20 calls it
thoracic only if one rib exceeds 38\,mm, and lumbar if both are
shorter~\citep{Liebl2021VerSeData}. The reference standard for numbering
remains a count from C2 on whole-spine imaging~\citep{Lian2018LSTVNumeration}. Reported LSTV
prevalence is 4--30\%~\citep{Konin2010LumbosacralRelevance},
21.1\% of 1{,}468 lumbar MRIs~\citep{Hanhivaara2022LSTV} and 16.3\% of 551 whole-spine CTs by
count~\citep{Nagata2025LSTVWholeSpine}. It varies with the population, from 8.1\% to 16.9\%
across three clinic groups in one series, which did not count Castellvi type~I as
transitional~\citep{Sekharappa2014LumbosacralSignificance}.

\paragraph{The six subtypes in these terms.}
The taxonomy of \S\ref{sec:dataset} inherits its sources' conventions. CTSpine1K numbers
C1--L6 as 1--25 and attributes an L6 to sacral lumbarization and a missing L5 to lumbar
sacralization~\citep{Deng2021CTSpine1K:Tomography}; \texttt{lumb} and
\texttt{sacr\_count} are these counts. A count does not place the anomaly: six lumbar
vertebrae can follow eleven thoracic ones, and four can follow a
T13 (\citealp{Moller2026VERIDAH:Sequences}, Figs.~1D and~5), and no scan here, all CT
colonography of the abdomen and pelvis, contains C2~\citep{Schwing2026CTSpinoPelvic1KArxiv}.
\texttt{sacralization} and \texttt{semisacralization} rest on CTPelvic1K's undocumented
filename tag. The April release has no thoracic, T13 or rib class, so the thoracolumbar border
cannot be checked on it; the current release (v11) adds thoracic levels, ribs, lumbar ribs and
a T13 class~\citep{Schwing2026CTSpinoPelvic1KArxiv}.

\paragraph{Castellvi classification.}
\citet{Castellvi1984LumbosacralDefects} typed the transitional segment mainly by the
morphology of its transverse process. As restated by \citet{Konin2010LumbosacralRelevance}, with \emph{a} marking a
unilateral and \emph{b} a bilateral finding, type~I (Ia, Ib) is a dysplastic transverse
process at least 19\,mm in craniocaudal width. Type~II (IIa, IIb) is an enlarged transverse process forming a
diarthrodial joint with the sacrum. Type~III (IIIa, IIIb) is complete osseous fusion of the
transverse process to the sacrum. Type~IV is type~II on one side and type~III on the other.
The grade ``does not provide information relevant to accurate enumeration of the involved
segment''~\citep[p.~1778]{Konin2010LumbosacralRelevance}, which is why \S\ref{sec:castellvi}
treats it as orthogonal to the count; type~IIIb occurs here at lumbar counts of four, five
and six (Table~\ref{tab:castellvi_dist}). VerSe'20 classified only types II--IV, lists each
subject's variants, and left vertebrae partly fused to the sacrum (III, IV)
unsegmented~\citep{Liebl2021VerSeData}; \VERIDAH{}'s reference labels
call any Castellvi-bearing vertebra a sacralized lumbar one, and it outputs no
grade~\citep{Moller2026VERIDAH:Sequences}.

\paragraph{\VERIDAH{} and SPINEPS.}
\VERIDAH{}~\citep{Moller2026VERIDAH:Sequences} labels vertebrae whose positions are given
(the paper takes them from instance segmentations). A 3-D DenseNet169 classifies a crop
centred on each vertebra with four heads: vertebra label (no T13 or L6 class), region,
transition (last cervical, first thoracic, last thoracic, first lumbar, last lumbar, or none)
and visibility. For training, T13 and L6 are relabelled T12 and L5, so an extra vertebra is a repeated T12 or L5 that only the transition
head separates, and anomalous cases are drawn four times per epoch. A sequence predictor
combines the heads into a path cost weighted by each vertebra's visibility and minimises it by
dynamic programming under hard constraints: consecutive, spatially ordered labels; 11--13
thoracic and 4--6 lumbar; at most one repeated T12 and one repeated L5, renamed T13 and L6.
On 365 in-house test CTs enriched for anomalies it labelled every vertebra correctly in
99.18\% of subjects, against 77.81\% for the method of \citet{Meng2022VertebraeMLMI}, and
recalled 97.22\% of lumbar enumeration anomalies; in its ablation one label head without the
predictor recalls 30.56\%, with it 63.89\%, and four heads without it 30.56\%.
SPINEPS~\citep{Moller2024SPINEPSautomaticSegmentation} is a semantic-then-instance
segmentation of T2-weighted MRI whose instances are counted top to bottom ``regardless of the
field of view or the presence of enumeration abnormalities''; \VERIDAH{} supplies the
labelling step and ships with its code~\citep{Moller2026VERIDAH:Sequences}. SPINEPS~2.1.2
(Appendix~\ref{app:spineps}) runs CT models before the \VERIDAH{} labeller and writes \VerSe{}
ids, L6 as 25 and T13 as 28~\citep{MollerSPINEPSSoftware}; the peer-reviewed SPINEPS paper
covers MRI only.

\paragraph{Why a missing class shifts labels.}
A segmenter gives each voxel one label from a fixed vocabulary. \VerSe{} numbers L1--L6 as
20--25 and T13 as 28~\citep{VerSeRepo}, and CTSpine1K and the April release have an L6 class;
the \texttt{total} task of \TS{} (version 2.4.0 was audited) has C1--L5, the sacrum and S1 but
no L6 or T13~\citep{Wasserthal2023TotalSegmentator:Images,TotalSegmentatorCode}.\footnote{Since
version 2.15.0 (1 July 2026) the \TS{} package has also defined a \texttt{total\_v3} task whose
class map puts \texttt{vertebrae\_L6} where \texttt{total} has \texttt{vertebrae\_S1}; its
weights were posted as a GitHub pre-release on 7 September 2026, and neither the changelog
nor the README describes it~\citep{TotalSegmentatorCode}. It postdates the
benchmark of \S\ref{sec:audit}; Appendix~\ref{app:tsv3} evaluates it under the same protocol.} Given
six lumbar-type segments and five lumbar labels, a model must reuse a label or give a segment
a thoracic, sacral or no label; on lumbarized spines \TS{} gives the L6 the L5 label in two of
three scorable junctions, and its lumbar Dice falls level by level toward the sacrum
(\S\ref{sec:audit}). An L6 output label is necessary but not sufficient, and it need not be a
trained class. CTSpine1K's own nnU-Net baseline scored an L6 Dice of almost
zero~\citep{Deng2021CTSpine1K:Tomography}; our unmerged control, trained with an L6
class, still labels most ground-truth L6 voxels in its validation patches L5
(\S\ref{sec:unmerged}); and \VERIDAH{} as published has no L6 class, yet names an L6 as a
repeated L5 that its transition head and sequence predictor
place~\citep{Moller2026VERIDAH:Sequences}.
